%webamc CODE=r203/2023/dhcp
%webamc TITLE=DHCP
%webamc RND=0

\noindent\begin{minipage}{.42\textwidth}
  On considère dans cet exercice le réseau de la figure ci-contre. Ce
  réseau est décomposé en deux sous-réseaux interconnectés par un
  routeur R. L'hôte relai est un relai DHCP pour son sous-réseau et
  l'hôte dhcp et un serveur DHCP pour tout le réseau. Les hôtes anna
  et eve sont configurées dynamiquement.
\end{minipage}
\begin{minipage}{.57\textwidth}
  \begin{center}\centering\begin{tikzpicture}[node distance=1.5cm]
  \switch{}{S1}{90}{switch\textsubscript{1}}
  \node[below of=S1] (pt) {};
  \machine{left of=pt}{anna}{270}{anna}
  \serveur{right of=pt}{relai}{270}{relai}
  \routeur{right of=S1,node distance=2.5cm}{R}{90}{R}
  \switch{right of=R,node distance=2.5cm}{S2}{90}{switch\textsubscript{2}}
  \node[below of=S2] (pt) {};
  \machine{left of=pt}{eve}{270}{eve}
  \serveur{right of=pt}{dhcp}{270}{dhcp}
  %
  \draw[liaison] (S1) -- node[pos=0.8,left] {\scriptsize eth0} (anna);
  \draw[liaison] (S1) -- node[pos=0.8,right] {\scriptsize eth0} (relai);
  \draw[liaison] (S1) -- node[pos=0.8,above] {\scriptsize eth0} (R);
  \draw[liaison] (S2) -- node[pos=0.8,above] {\scriptsize eth1} (R);
  \draw[liaison] (S2) -- node[pos=0.8,left] {\scriptsize eth0} (eve);
  \draw[liaison] (S2) -- node[pos=0.8,right] {\scriptsize eth0} (dhcp);
\end{tikzpicture}
\end{center}
\end{minipage}

Anna demande un bail. On capture alors les messages DHCP
ci-dessous. Seul le début de l'échange est représenté.
\begin{center}\begin{tempDiag}[>=stealth'][start=1,numbers=1]{5}{3cm}
  \tempDiagNode{1}{10pt}{relai}{relai}
  \tempDiagNode{2}{10pt}{anna}{anna}
  \tempDiagNode{3}{10pt}{routeur}{R}
  \tempDiagNode{4}{10pt}{dhcp}{dhcp}
  \tempDiagNode{5}{10pt}{eve}{eve}
  %
  \tempDiagMsg[shift=10,number=0]{2}{1,3}{}
  \node[right of=pt-2] (pt) {};
  \node[above of=pt,number,node distance=5mm] {1};
  \tempDiagMsg{1}{3,4}{}
  \tempDiagMsg{4}{3,1}{}
  \tempDiagMsg{1}{2}{}
\end{tempDiag}
\end{center}

Voici les options (en hexadécimal) inclues dans le message 1:
\begin{center}\tt 35 01 01 33 01 06 37 03 2a 0f 01 ff\end{center}

Voici les options (en hexadécimal) inclues dans le message 3:
\begin{center}\tt 35 01 02 33 01 03 36 04 0f 01 00 20 01 04 ff ff f0 00 2a 04 0f 01 06 2a ff\end{center}
                  
